
\section{Sujets type bac}


\begin{exoBox}{bleu}{Sujet 1 -- Comparaison de deux fonctions exponentielles}
On considère les fonctions $f$ et $g$ définies sur $\R$ par :
\[
f(x)=e^{2x}+e^2
\qquad\text{et}\qquad
g(x)=(1+e^2)e^x.
\]
On note $\mathcal C_f$ et $\mathcal C_g$ leurs courbes représentatives.
\begin{enumerate}
\item Montrer que, pour tout réel $x$ :
\[
f(x)-g(x)=e^{2x}-(1+e^2)e^x+e^2.
\]
\item En posant $X=e^x$, factoriser le trinôme :
\[
X^2-(1+e^2)X+e^2.
\]
\item En déduire que, pour tout réel $x$ :
\[
f(x)-g(x)=(e^x-1)(e^x-e^2).
\]
\item Déterminer les abscisses des points d'intersection de $\mathcal C_f$ et $\mathcal C_g$.
\item Donner les coordonnées exactes de ces points d'intersection.
\item Étudier le signe de $f(x)-g(x)$ selon les valeurs de $x$.
\item En déduire la position relative de $\mathcal C_f$ et $\mathcal C_g$.
\item Résoudre l'inéquation $f(x)\le g(x)$.
\item Sans calculer de valeurs décimales, comparer $2e^2$ et $e+e^3$.
\end{enumerate}
\end{exoBox}

\begin{exoBox}{vert}{Sujet 2 -- Étude d'une fonction avec exponentielle}
On considère la fonction $f$ définie sur $[-2;4]$ par :
\[
f(x)=(x-2)e^x.
\]
\begin{enumerate}
\item Calculer $f(0)$.
\item Montrer que, pour tout réel $x$ :
\[
f'(x)=(x-1)e^x.
\]
\item Étudier le signe de $f'(x)$ sur $[-2;4]$.
\item Dresser le tableau de variations de $f$ sur $[-2;4]$.
\item Déterminer la valeur exacte du minimum de $f$ sur $[-2;4]$.
\item Étudier le signe de $f(x)$ sur $[-2;4]$.
\item Déterminer l'équation de la tangente à la courbe de $f$ au point d'abscisse $0$.
\end{enumerate}
\end{exoBox}

\begin{exoBox}{orange}{Sujet 3 -- Modèle en contexte}
La concentration d'un produit dans le sang, en unité arbitraire, est modélisée pendant les cinq premières heures par la fonction $C$ définie sur $[0;5]$ par :
\[
C(t)=t e^{1-t}.
\]
\begin{enumerate}
\item Calculer $C(0)$ et $C(1)$.
\item Montrer que, pour tout $t\in[0;5]$ :
\[
C'(t)=(1-t)e^{1-t}.
\]
\item Étudier le signe de $C'(t)$ sur $[0;5]$.
\item Dresser le tableau de variations de $C$ sur $[0;5]$.
\item À quel instant la concentration est-elle maximale ? Quelle est cette concentration maximale ?
\item Calculer $C(2)$.
\item On admet que $e<3$. Montrer que $C(2)>\dfrac23$.
\item Interpréter ce résultat dans le contexte.
\end{enumerate}
\end{exoBox}
